\documentclass[10pt,a4paper]{amsart}
\usepackage[margin=19mm]{geometry}
\usepackage{amsmath,amssymb,mathtools}
\usepackage[scr=rsfs,cal=euler]{mathalfa}
\usepackage{xfrac}
\usepackage[unicode,hidelinks,bookmarks=false]{hyperref}
\newcommand{\ie}{i.e.\ }
\newcommand{\cf}{cf.\ }
\newcommand{\resp}{respectively\ }
\newcommand{\Subsection}{\subsection}
\providecommand{\nobreakdash}{\nobreak\hbox{-}}
\setlength{\emergencystretch}{2em}
\multlinegap0pt
\usepackage{siunitx}
\usepackage{siunitx}
\usepackage{siunitx}
\usepackage{siunitx}
\usepackage{amssymb}
\usepackage{mathtools}
\usepackage{siunitx}
\usepackage[shortlabels]{enumitem}
\usepackage{algorithm}\usepackage{algpseudocode}
\newtheorem{definition}{Definition}
\newtheorem{theorem}{Theorem}
\newcommand{\bs}[1]{\boldsymbol{#1}}
\title{Explicit Construction of Approximate Kolmogorov Superpositions with C2 Smoothness}
\author{Lunji Song$^\dagger$, Zilan Cheng$^\ddagger$, Juan Diego Toscano$^\S$, Li-Lian Wang$^\ddagger$}
\date{Source: arXiv:2508.04392v2 (CC BY 4.0)}
\begin{document}
\maketitle

\noindent\emph{Source and licence.} Explicit Construction of Approximate Kolmogorov Superpositions with C2 Smoothness. Version arXiv:2508.04392v2 (CC BY 4.0), distributed by arXiv under the Creative Commons Attribution 4.0 International licence, http://creativecommons.org/licenses/by/4.0/, as recorded in the arXiv licence metadata for this exact version and shown on the abstract page https://arxiv.org/abs/2508.04392v2. Full licence notice: \texttt{licenses/arxiv-2508.04392-CC-BY-4.0.txt}.\par\smallskip

\noindent\textit{Editorial scope.} Counted unit: the article's theorem MainResult (source0.tex lines 1329--1340). The article's complete original proof of it is delivered with it (source0.tex lines 1341--1420, one \texttt{proof} environment of 80 lines). No mathematics is written, repaired or re-derived: the statement and the proof are the article's own lines, in the article's own notation, and no retained line is substituted or re-flowed.  Retained uncounted as the source-local context the proof needs: lines 344--346; lines 386--399; lines 468--480; lines 586--587; lines 592--594; lines 750--752; lines 1022--1031; lines 1108--1111; lines 1117--1132; lines 1265--1279.  Nothing was omitted from any retained span, so the package carries no omission record.  No retained line contains a citation, so the wrapper carries no bibliography and no bibliography entry.  No retained line contains a figure environment, an image-inclusion command or drawing code, so no image is delivered and the package declares no asset dependency.  Editorial formatting only, in addition: an AMS \texttt{amsart} preamble that restates the article's own declarations and macros used by the retained lines, namely the environments \texttt{definition}, \texttt{theorem} on separate counters, together with the standard packages those lines need.  The retained environments print in this document as Definition 1, Theorem 1 in order of appearance, whereas the article prints the same items with the numbering of its own full document.  The complete native source of the article as distributed in the arXiv e-print for this exact version is bundled unchanged as \texttt{sources/182799/source0.tex}.
\begin{equation}\label{St}
    S(z):=z-\frac{1}{2\pi} \sin{2\pi z},\quad z\in [0,1].
\end{equation}

\begin{equation}\label{shapfunc0}
\begin{split}
\phi(t)\big|_{t\in B^j}&=
\begin{dcases}
y_j+(y_j^\nu-y_j)\, S\big(\tfrac{t-t_j}{t_j^\delta-t_j}\big), &t\in G^j,\\
y_j^\nu+(y_{j+1}-y_j^\nu)\,P\big(\tfrac{t-t_j^\delta}{t_{j+1}-t_j^\delta}\big), &t\in I^j,
\end{dcases}\\[4pt]
&=
\begin{dcases}
jh+\nu\, S\big(\tfrac{t-jh}\delta\big), &t\in [jh, jh+\delta],\\
jh+\nu+(h-\nu)\,P\big(\tfrac{t-jh-\delta}{h-\delta}\big), &t\in [jh+\delta, (j+1)h].
\end{dcases}
\end{split}
\end{equation}

\begin{definition}[{\bf Inner functions}]\label{Inner-fun}  \emph{Let $\phi(t)$ be the piecewise $C^2,$ strictly increasing function defined in
\eqref{shapfunc0}, and introduce the linear transformations: \begin{equation}\label{shift-coord}
x=t+(q-1)\delta\;\;\; {\rm or} \;\;\; t=x+(1-q)\delta, \quad  t\in B^j,\;\; x\in B_q^j.
\end{equation} 
We define the inner function  at the $q$-th level by
\begin{equation}\label{PsiqxAS}
\begin{split}
\psi_q(x)& =\frac{\phi(x+(1-q)\delta)- \phi((1-q)\delta)} {\phi(1+(1-q)\delta)-\phi((1-q)\delta)}\\[4pt]
&=\phi(x+(1-q)\delta)- \phi((1-q)\delta), \quad  x\in B_q^j\cap [0,1],
\end{split}
\end{equation}
 for $j\in \mathcal J$ and $q=1,\ldots, N.$} 
\end{definition}

\begin{equation}\label{cqjcenter}  \bs c_q^{\bs j}:=(c_q^{j_1},\ldots, c_q^{j_d})\;\;\; {\rm with} \;\;\; c_q^{j_p}=j_ph+(h+\delta)/2+(q-1)\delta.
    \end{equation} 

\begin{equation}\label{lower-corner1}
   {\bs v}_{q}^{\bs j}=( v_q^{j_1},\ldots,  v_q^{j_d})=\bs j h+q\,\delta\,\bs 1,\quad \tilde{\bs v}_{q}^{\bs j}=(\tilde v_q^{j_1},\ldots, \tilde v_q^{j_d})=\bs v_q^{\bs j}+(h-\delta)\,\bs 1,
\end{equation}

\begin{equation}\label{mon-hyperb}
e_q^{\bs j}=\Psi_q(\bs v_q^{\bs j};\bs \lambda)  \le \Psi_q(\bs x;\bs \lambda) \le  \Psi_q(\tilde{\bs v}_q^{\bs j};\bs \lambda)=\tilde e_q^{\bs j},\quad \forall\, \bs x\in \bs C_q^{\bs j}.
\end{equation}

\begin{equation}\label{secondcondforanydd}
\begin{dcases}
 e_q^{(-1,j_2+1,j_3,\ldots,j_d)}-\tilde e_q^{(N-1,j_2,j_3\ldots,j_d)}
>-\lambda_1+\lambda_2\nu+(\lambda_3+\cdots+\lambda_d)(\nu-h),
  \\[6pt]
e_q^{(-1,-1,j_3+1,\ldots,j_d)}-\tilde e_q^{(N-1,N-1,j_3\ldots,j_d)}>-(\lambda_1+\lambda_2)+\lambda_3\nu+(\lambda_4+\cdots+\lambda_d)(\nu-h), \\
\qquad \qquad \vdots\\
e_q^{(-1,-1,\ldots,-1,j_d+1)}-\tilde e_q^{(N-1,N-1,\ldots,N-1,j_d)}>-(\lambda_1+\lambda_2+\ldots+\lambda_{d-1})+\lambda_d\nu.
\end{dcases}
\end{equation}

\begin{equation}\label{qAeqnA} 
 \nu=\nu^*:=-\frac{1-h}{2}+ \frac{1-h}{2}\sqrt{1+ \frac 
    {4h}{(1-h)^2}\Big(1 - \Big(\frac{h}{h+1}\Big)^{d-1}\Big)}\,,
\end{equation}

\begin{theorem}\label{Hypercube-Sep} For $d\ge 2$ and   $0<h<1,$ we take $\nu=\nu^*$ in \eqref{qAeqnA}, and set 
\begin{equation}\label{Spcurrent}
{\mathcal S}_1^*:=1,\quad {\mathcal S}_{p+1}^*=\frac{1}{2}\Big(Q^*_{p+1}+\frac{\nu^*+1}{h}{\mathcal S}_p^*-\frac{1}{h}\Big),\quad 1\le p\le d-1,
\end{equation}
with $Q_d^*=\tfrac 1 {\nu^*+1}.$ Letting 
\begin{equation}\label{lmdpS}
    \lambda_p^*={\mathcal S}_p^*-{\mathcal S}_{p+1}^*,\quad 1\le p\le d-1;\quad \lambda_d^*={\mathcal S}_d^*,
\end{equation}
we have $\bs \lambda^*\in {\mathbb A}_h(\nu^*),$ and 
\begin{equation}\label{def_d_lambda}
    \lambda_p^*=
        \frac{h^{d-p}}{(\nu^*+1)^{d-p+1}}+\chi_p\frac{\nu^*-h+(\frac{h}{\nu^*+1})^d}{\nu^*+1-h}\Big(\frac{\nu^*+1}{2h}\Big)^{p-1},
\end{equation}
where $\chi_p=\tfrac{2h-\nu^*-1}{2h}$ for $1\le p\le d-1$ and $\chi_p=1$ for $p=d$.
With this choice, the corresponding sequence $\bs E_q$ is strictly increasing. 
\end{theorem}

\begin{definition}[{\bf  Approximate KST}]\label{Defn:Outer01} {\em  For $d\ge 2$ and   $0<h<1,$
let $\nu_*\in (0,h)$ and $\bs \lambda^*=(\lambda_1^*, \ldots, \lambda_d^*)$  be  given in Theorem \ref{Hypercube-Sep}, and let $\{(\tau_k, s_k)\}_{k=1}^{K}$ be the elements of 
$(\bs A_q, \bs B_q)$ with $K$ being the length of $\bs A_q.$ Then we define the outer function at each level $1\le q\le N$ as 
\begin{equation}\label{qj2g01-A}
  g_q(t)=s_k +(s_{k+1}-s_k) S\big(\tfrac{t-\tau_k}{\tau_{k+1}-\tau_k}\big), \quad t\in  [\tau_k, \tau_{k+1}],
\end{equation}
for $k=1,\ldots, K-1,$ where $S$ is the smeared-out Heaviside function  defined in \eqref{St}. Accordingly, the approximate Kolmogorov  representation of $f\in C([0,1]^d)$  takes the form 
\begin{equation}\label{0fx-KA-form-B00}
      f_N(\bs x)=f_N(\bs x; \bs \lambda^*)=\sum_{q=1}^{N} g_q\circ\Psi_q(\bs x;\bs \lambda^*)\;\;\;
      {\rm with}
      \;\;\; \Psi_q(\bs x;\bs \lambda^*)=\sum_{p=1}^d \lambda_p^*\, \psi_q(x_p),
\end{equation}
where the inner functions $\psi_q$ are defined in Definition \ref{Inner-fun}.
}
\end{definition}

\begin{theorem}\label{MainResult} Let $d\ge 2$ and 
assume that  $f\in {\mathcal H}^\alpha([0,1]^d)$ is any $\alpha$-H\"older continuous function satisfying 
$f\in C([0,1]^d)$ and 
\begin{equation}\label{alpha-cond}
    |f(\bs x)-f(\bs y)|\le c_{\alpha}\, \|\bs x-\bs y\|_2^\alpha, \quad \forall\, \bs x,\bs y \in [0,1]^d,\;\;\; 0<\alpha\le 1.
\end{equation}
 Let $f_N(\bs x)$ be the constructed superpositions as in Definition \ref{Defn:Outer01}. Then we have 
\begin{equation}\label{approximA}
\|f-f_N\|_\infty \le  \Big(1+\frac{d^{\frac{\alpha}{2}}}{2^\alpha} +\frac{2 d^{\frac \alpha 2+1 }}{N^{1-\alpha}}\Big)\frac {c_{\alpha}} {N^{\alpha}}.
\end{equation}
Here, $\|\bs z\|_2$ is the length of the vector $\bs z$, and $\|\cdot\|_\infty$ is the usual $L^\infty$-norm.  
\end{theorem}

\begin{proof} 
By definition,  each $\bs x\in [0,1]^d$ is covered by a hypercube   
$\bs C_q^{\bs j}$ for  at least $(N-d)$ different values of $q\in\{1,\ldots, N\}$, but not covered by at most $d$ values of $q$ (i.e., falls into a gap). 
Then we can split the error into 
\begin{equation}\label{ffNas}
\begin{aligned}
|f(\bs x)-f_N(\bs x)|& = \Big|f(\bs{x})-\sum_{q=1}^N g_q\circ\Psi_q(\bs{x})\Big|\\
& =
\Big|\sum_{q=1}^N \Big(\frac {f(\bs{x})} N- g_q\circ\Psi_q(\bs{x})\Big)\Big| \leq \sum_{q=1}^N\Big|\frac{f(\bs x)}{N}-g_q\circ\Psi_q(\bs x)\Big| \\
& =  \sum_{q\in {\mathcal Q}}\Big|\frac{f(\bs x)}{N}-g_q\circ\Psi_q(\bs x)\Big|+\sum_{q\not\in {\mathcal Q}}\Big|\frac{f(\bs x)}{N}-g_q\circ\Psi_q(\bs x)\Big|,
\end{aligned}
\end{equation}
where  ${\mathcal Q}$ denotes the set of at least $N-d$ values of $q$ for which the hypercubes $\bs C_q^{\bs{j}}$ cover the point $\bs x,$ and the second summation over $q\not\in {\mathcal Q}$ counts the contributions of the levels where $\bs x$ is in a gap.  

If $q\in {\mathcal Q},$ i.e.,  $\bs x\in  \bs C_q^{\bs{j}}$  for some $\bs j$ with the center $\bs{c}_{q}^{\bs{j}}$ (see \eqref{cqjcenter}),  it follows from \eqref{alpha-cond}  that
\begin{equation}\label{Oscf2}
|f(\bs{x})-f(\bs{c}_q^{\bs{j}})|\leq c_{\alpha}\,\|\bs{x}-\bs{c}_q^{\bs{j}}\|^{\alpha}_2 \leq c_{\alpha} \Big(\frac{\sqrt{d}(N-1)\delta}{2}\Big)^{\alpha}.
\end{equation}
Correspondingly, we can identify from  \eqref{mon-hyperb} that the image $t=\Psi_q(\bs{x};\bs \lambda^*)\in [e_q^{\bs j}, \tilde e_q^{\bs j}]$ must lie in a piece, say $\tau_k, \tau_{k+1}$ in \eqref{qj2g01-A}.  Thus,   
 we  derive from   \eqref{Oscf2} that 
\begin{equation}\label{Nsquare2}
\begin{split}
\Big|\frac{f(\bs{x})}{N}  -g_q\circ\Psi_q(\bs{x})\Big| & \le  \Big|\frac {f(\bs{x})}N - \frac   {f(\bs{c}_q^{\bs{j}})} N\Big|+\Big|\frac   {f(\bs{c}_q^{\bs{j}})} N-g_q\circ\Psi_q(\bs{x};\bs \lambda^*)\Big| \le \frac 1 N \big|{f(\bs{x})}-   {f(\bs{c}_q^{\bs{j}})} \big|\\[4pt]
&\quad +\big|(s_{k+1}-s_k) S\big((\Psi_q(\bs x;\bs \lambda^*)-\tau_k)/(\tau_{k+1}-\tau_k)\big)\big|\\[4pt]
&\le \frac 1 N \big|{f(\bs{x})}-   {f(\bs{c}_q^{\bs{j}})} \big|+|s_{k+1}-s_k | \\[4pt]
&\le  \frac{c_{\alpha}} N \Big(\frac{\sqrt{d}(N-1)\delta}{2}\Big)^{\alpha}+ \frac{c_{\alpha}}{N} (N\delta)^{\alpha}
\\[4pt]
&\leq \frac{c_{\alpha}}{N}  \Big(1+\frac{d^{\frac{\alpha}{2}}} {2^{\alpha}}\Big)(N\delta)^{\alpha},
\end{split}
\end{equation}
where we have used $S(z)\in [0,1]$ (noted: $|\bs \lambda^*|_1=1$) and  \eqref{alpha-cond} with the function values in the corresponding row:  
$$s_{k+i}=b_{q,k+i}^{\bs j}=\tfrac{f(\bs{c}_{q,k+i}^{\bs{j}})}{N},\quad i=1,2,$$  to bound
\begin{equation}\label{bqjest}
\begin{split}
|s_{k+1}-s_k | &=|b_{q,k+1}^{\bs j,j_p}-b_{q,k}^{\bs j,j_p}|=\frac 1 N\big|f(\bs c_{q,k+1}^{\bs j,j_p})-f(\bs c_{q,k}^{\bs j,j_p})\big|\\[4pt]
  &\le \frac {c_{\alpha}} N \big\|\bs c_{q,k+1}^{\bs j,j_p}-\bs c_{q,k}^{\bs j,j_p}\big\|_2^\alpha  \leq \frac{c_{\alpha}}{N} (N\delta)^{\alpha}.
  \end{split}
\end{equation}
Then we obtain estimate of the first term in \eqref{ffNas}:
\begin{equation}\label{term1-44}
    \sum_{q\in {\mathcal Q}}\Big|\frac{f(\bs x)}{N}-g_q\circ\Psi_q(\bs x)\Big| \le \frac{c_{\alpha}(N-d)}{N}  \Big(1+\frac{d^{\frac{\alpha}{2}}} {2^{\alpha}}\Big)(N\delta)^{\alpha}.
\end{equation}
It remains to consider  $q\notin {\mathcal Q},$ i.e.,
$\bs x$ is in a gap. Again letting  
$t=\Psi_q(\bs{x};\bs \lambda^*),$ 
we can identify the corresponding piece $t\in [\tau_\ell,\,  \tau_{\ell+1}],$ associated with the centers and data:
$$
\tau_\ell= \Psi_q(\bs{c}_q^{\bs i};\bs \lambda^*),\quad  \tau_{\ell+1}= \Psi_q(\bs{c}_q^{\bs i'};\bs \lambda^*), \quad  s_\ell= f(\bs{c}_q^{\bs i})/N,\quad  s_{\ell+1}= f(\bs{c}_q^{\bs i'})/N.
$$
We choose the center closest to $\bs x,$ say $\bs{c}_q^{\bs i}.$ Following the first three lines of \eqref{Nsquare2}, we have 
\begin{equation}\label{Nsquare2-00}
\begin{split}
\Big|\frac{f(\bs{x})}{N}  -g_q\circ\Psi_q(\bs{x})\Big| & 
\le \frac 1 N \big|{f(\bs{x})}-   {f(\bs{c}_q^{\bs{i}})} \big|+|s_{\ell+1}-s_\ell|\\[4pt] 
& =\frac 1 N \big|{f(\bs{x})}-   {f(\bs{c}_q^{\bs{i}})} \big|+\frac 1 N \big|f(\bs{c}_q^{\bs i})-   f(\bs{c}_q^{\bs{i'}}) \big|\\[4pt]
&\le \frac{c_{\alpha}} N \big(\|\bs x- 
\bs{c}_q^{\bs{i}}\|_2^\alpha +
\|\bs{c}_q^{\bs{i}}-\bs{c}_q^{\bs{i'}}\|_2^\alpha\big).
\end{split}
\end{equation}
Note that the distance between $\bs{c}_q^{\bs{i}}$ and $\bs{c}_q^{\bs{i'}}$ can be controlled by the longest distance between two adjacent $\bs v_q^{\bs j}$ and 
$\tilde {\bs v}_q^{\bs j}$ in 
\eqref{lower-corner1}. From \eqref{secondcondforanydd}, we find  
\begin{equation}\label{cqi-est}
\|\bs{c}_q^{\bs{i}}-\bs{c}_q^{\bs{i'}}\|_2
\le \| \bs v_q^{(-1,-1,\ldots,-1,j_d+1)}-\tilde {\bs v}_q^{(N-1,N-1,\ldots,N-1,j_d)} \|_2
\le \sqrt{d-1+\delta^2}.
\end{equation}
The choice of $\lambda^*$ can ensure the ordering of images of the points in the hypercubes in rows, but cannot guarantee the points in the gaps. In other words, although $t\in [\tau_\ell,\,  \tau_{\ell+1}],$  its pre-image $\bs x$ may be far away from the pre-images $\bs{c}_q^{\bs{i}}$ and $\bs{c}_q^{\bs{i'}}$ (e.g, for the gaps of the largest volume  $(h-\delta)^{d-1}\delta$), the worst case can always be bounded by 
\begin{equation}\label{cqi-est2}
\|\bs x- 
\bs{c}_q^{\bs{i}}\|_2\le \sqrt d.
\end{equation}
Thus from  \eqref{Nsquare2-00} and the above,  we can obtain  
\begin{equation}\label{term1-44-00}
    \sum_{q\not \in {\mathcal Q}}\Big|\frac{f(\bs x)}{N}-g_q\circ\Psi_q(\bs x)\Big| \le  \frac{2c_{\alpha}}{N} d^{\frac \alpha 2+1}.  
\end{equation}
Then the estimate \eqref{approximA} is a direct consequence of  \eqref{ffNas}
\eqref{alpha-cond}, \eqref{term1-44} and \eqref{term1-44-00}. 
\end{proof}

\end{document}
